\bookchapter{الفصل 35: الأطفال والهاتف — دليل التربية في العصر الرقمي}

\chapterintro{هذا الفصل لا يقول “الهاتف وحش” ولا يروج “التعرف المبكر يزيد الذكاء”. نقسم حسب العمر: إعطاء أجهزة، وضع قواعد، تعليم آداب، حماية صحة، خطوة بخطوة.}

\booksection{أولاً: متى نعطي الطفل هاتفاً؟}

لا إجابة قياسية، لكن هناك خط أساس عمري. إعطاء الهاتف مبكراً، قشرة الدماغ الجبهية للطفل لم تنضج، ضبط النفس لا يواكب “تصميم الإغراء” للخوارزميات؛ والتأخير قد يجعله “منقطع اجتماعياً” بين أقرانه. الاقتراحات العمرية:

\begin{itemize}
  \item \textbf{3-6 سنوات}: تجنب لمس الهاتف قدر الإمكان. إن لزم، لا تتجاوز الجلسة 15 دقيقة، المحتوى رسوم تعليمية منتقاة من الوالدين، ويجب\textbf{المشاهدة المرافقة}. هذه المرحلة التأثير على الرؤية والانتباه أكبر.
  \item \textbf{7-9 سنوات}: يمكن "الاستعارة" على جهاز الوالدين، للدراسة عبر الإنترنت والاستماع للقصص والبحث. ابدأ بقاعدة "اطلب قبل الاستخدام وأعد بعد الانتهاء"، لا تخصص هاتفاً خاصاً.
  \item \textbf{10-12 سنة}: فكّر في إعطاء "هاتف وظائف" أو ساعة طفل، للتواصل ذهاباً وإياباً للمدرسة. إن لزم هاتف ذكي، فعّل رقابة الوالدين الصارمة.
  \item \textbf{13 سنة فأكثر}: دخول الإعدادي، الحاجة الاجتماعية تزداد، يمكن إعطاء هاتف ذكي، لكن مع "ميثاق استخدام" ومراجعة دورية.
\end{itemize}

\begin{tcbinfo}{مبدأ "يكفي ما يكفي"}
هاتف الطفل الأول لا يحتاج أن يكون رائد الأعمال. “الأداء الفائض” للمعالج غالباً يُستخدم للألعاب. اختر هاتفاً ببطارية كبيرة وأداء متوسط وشاشة مريحة للعين في فئة الألف أو جهاز مُجدّد رسمي، أفضل للتحكم. الميزانية: 1500-2500 يوان تكفي أغلب الطلاب.

\end{tcbinfo}

\booksubsection{الساعة الذكية كحل انتقالي}

إذا كنت تريد حل “هل يمكن التواصل بعد المدرسة”، ساعة الطفل أنسب من الهاتف. تتصل وترسل صوتاً وتحدد موقعاً، دون متصفح وفيديوهات قصيرة، تحد “مداخل الإدمان” طبيعياً:

\begin{itemize}
  \item \textbf{ساعة XiaoTiancai}: منظومة أكمل، "المس لإضافة صديق" ميزتها الأساسية، لكن منظومة التطبيقات مغلقة نسبياً.
  \item \textbf{ساعة Huawei للأطفال}: ترابط جيد مع هواتف Huawei، تحديد موقع دقيق، تدعم مكالمات فيديو، قيمة جيدة.
  \item \textbf{ساعة Xiaomi للأطفال}: سعر مناسب، وظائف أساسية كاملة، مناسبة للميزانيات المحدودة.
\end{itemize}

عند الشراء انتبه: اختر موديلاً يدعم\textbf{4G جميع الشبكات}(2G/3G تتقاعد تدريجياً)، وتأكد من باقة الطفل من المشغل.

\booksection{ثانياً: دليل شامل لنظام رقابة الوالدين}

قبل إعطاء الجهاز للطفل، أول شيء ليس التشغيل، بل ضبط رقابة الوالدين. لكل من المنصتين حل أصلي:

\booksubsection{iOS: وقت استخدام الشاشة}

\begin{itemize}
  \item \textbf{حد التطبيقات}: الإعدادات وقت استخدام الشاشة حد التطبيقات، يمكن ضبط مدة يومية لكل فئة (مثل ألعاب 30 دقيقة).
  \item \textbf{وقت الإيقاف}: اضبط "فترة إطفاء" يومية (مثل 22:00-7:00)، خلالها تُعطّل جميع التطبيقات، الهاتف فقط متاح.
  \item \textbf{قيود المحتوى}: تصفية الأفلام والصفحات وتثبيت التطبيقات حسب العمر، ومنع الشراء داخل التطبيقات.
  \item \textbf{المشاركة العائلية}: استخدم iPhone الوالد لإدارة جهاز الطفل عن بعد، الموافقة على طلبات التحميل وعرض تقارير الاستخدام.
\end{itemize}

\booksubsection{Android: Google Family Link / رقابة الوالدين للعلامة}

\begin{itemize}
  \item \textbf{Google Family Link}: حل Google الرسمي، إدارة مدة الاستخدام عن بعد وقفل الجهاز وعرض تقارير النشاط. Android المحلي حسب التوافر.
  \item \textbf{مساحة Xiaomi للأطفال}: الإعدادات مساحة الأطفال، ضبط مدة الاستخدام وتطبيقات القائمة البيضاء، الخروج بكلمة مرور الوالد.
  \item \textbf{وضع Huawei للطلاب}: الإعدادات الاستخدام الصحي للهاتف وضع الطالب، تقييد فترات الاستخدام والتطبيقات، يدعم الإدارة عن بعد.
  \item \textbf{وضع OPPO/vivo للأطفال}: يُفعل من "الإعدادات" أو تطبيق "مساحة الأطفال" المخصص، المبدأ مشابه.
\end{itemize}

\begin{center}
\begin{tabularx}{\linewidth}{|X|X|X|}
\hline
\textbf{الوظيفة} & \textbf{iOS وقت الشاشة} & \textbf{Android Family Link / حلول العلامات} \tabularnewline
\hline
حد التطبيقات & مدعوم، حسب الفئة & مدعوم، حسب التطبيق \tabularnewline
وقت الإيقاف & مدعوم & مدعوم (حلول العلامات تتفاوت) \tabularnewline
تصفية المحتوى & على مستوى النظام، أكثر شمولاً & يعتمد على العلامة، متفاوت \tabularnewline
إدارة عن بعد & المشاركة العائلية & Family Link / سحابة العلامة \tabularnewline
موافقة تثبيت التطبيقات & تتطلب موافقة الوالد & Family Link يدعم \tabularnewline
\hline
\end{tabularx}
\end{center}

\begin{tcbwarn}{رقابة الوالدين ليست "مرة واحدة للأبد"}
الطفل يجد دائماً طريقة للتجاوز (ابحث “كيف تكسر وقت الشاشة” ستجد دروساً كثيرة). الوسائل التقنية مجرد حد أدنى، الدفاع الحقيقي هو تواصل الوالدين والطفل والثقة. راجع التقرير شهرياً مع الطفل، بدلاً من المراقبة الأحادية.

\end{tcbwarn}

\booksection{ثالثاً: وضع الشباب}

منصات المحتوى الرئيسية تدمج “وضع الشباب”، عند التفعيل يحدد مدة الاستخدام ويصفي المحتوى السيئ ويغلق الدفع والإكرام. هذه وظيفة إلزامية بالقانون، لكن\textbf{تحتاج الوالد لتفعيلها وضبط كلمة المرور}:

\begin{itemize}
  \item \textbf{وضع Douyin/Kuaishou للشباب}: 40 دقيقة يومياً (يُحظر 22:00-6:00)، المحتوى فيديوهات علم وتعليم مختارة فقط، لا بث مباشر ولا إكرام ولا شحن.
  \item \textbf{وضع WeChat للشباب}: تقييد WeChat Pay (يلزم تفويض الوالد لمبلغ)، تقييد محتوى الحسابات العامة وقنوات الفيديو، هزّة و"الأشخاص القريبون" تُغلق افتراضياً.
  \item \textbf{نظام منع إدمان الألعاب}: مصادقة هوية على المستوى الوطني+تحديد وقت. القاصرون يمكنهم تسجيل الدخول فقط الجمعة والسبت والأحد والعطل 20:00-21:00، ساعة يومياً.
\end{itemize}

\textbf{لا تستخدم هوية الوالد للمصادقة “استعارة حساب”}\par

كثير من الأطفال يستخدمون هوية الوالد لتسجيل حساب لعبة لتجاوز منع الإدمان، هذا يخالف الاتفاقية ويُفقد الطفل الحماية. احتفظ بكلمة مرور الدفع ومعلومات الهوية، لا “تستعير حساباً” للطفل للسهولة.

\booksection{رابعاً: تعليم آداب الإنترنت}

التقنية تحجب 90\% من المخاطر، و10\% الباقية تعتمد على حكم الطفل. بدل “لا تتصفح”، الأهم تعليمه “كيف يتصفح”:

\begin{itemize}
  \item \textbf{حماية الخصوصية}: لا تكشف لغريب اسماً حقيقياً أو مدرسة أو عنوان منزل؛ لا ترسل صوراً شخصية أو هوية أو صور عائلة؛ تحديد الموقع مغلق افتراضياً، يُفعل عند الحاجة فقط.
  \item \textbf{التعرف على التنمر الإلكتروني والتعامل معه}: الشتم والعزل والتشويه بالصور كلها تنمر. أخبر الطفل عند مواجهتها لا ترد بمثل، لا تحذف الأدلة، أخبر الوالد فوراً. في WeChat اضغط مطولاً على الرسالة للشكوى، سيتعامل المنصة مع المتنمر.
  \item \textbf{نقد المحتوى}: علّم الطفل التمييز بين "الحقيقة" و"الرأي"، التعرف على العناوين المثيرة والصادمة؛ فهم أن الخوارزمية تعرض ما تريد رؤيته، لا يمثل العالم كله؛ أي إعلان "هدية من السماء" لا يُصدَّق.
  \item \textbf{الأمن السيبراني}: لا تنقر روابط غريبة، لا تمسح رموز QR مجهولة، لا تثبّت تطبيقات من خارج المتجر الرسمي؛ كلمة المرور لا تُشارك مع أحد (حتى "الصديق المقرب")؛ من يطلب رمز التحقق يُعتبر محتالاً.
\end{itemize}

\begin{tcbpractice}{جرّب: وقّع مع الطفل "ميثاق استخدام الهاتف"}
كتابة القواعد وتوقيع الطرفين ولصقها في مكان ظاهر، أكثر فاعلية من الاتفاق الشفهي. يُقترح أن يتضمن:

\begin{itemize}
  \item فترات الاستخدام: مثل "بعد الواجبات 19:30-20:30".
  \item مناطق المنع: على الطاولة، قبل النوم في غرفة النوم، في الفصل يُخبأ الهاتف.
  \item حدود المحتوى: لا تثبّت تطبيقات دون إخبار الوالد، لا تدردش مع الغرباء.
  \item عواقب المخالفة: مثل "تجاوز الوقت مرة، يُخفض وقت اليوم التالي للنصف".
  \item وعد الوالد: لا "يتصفح الهاتف" أمام الطفل فوق المدة المتفق عليها، قدوة.
\end{itemize}

\end{tcbpractice}

\booksection{خامساً: عادات الاستخدام الصحي}

عين الطفل ورقبته ونومه في طور النمو، ضرر الهاتف على جسمه أكبر من البالغين:

\begin{itemize}
  \item \textbf{اقتراح وقت الشاشة}: 3-6 سنوات لا تتجاوز ساعة يومياً؛ 6-12 سنة التعليم بلا حد، الترفيه لا يتجاوز ساعة؛ فوق 12 سنة تتفق مع الطفل، لكن الترفيه لا يتجاوز ساعتين.
  \item \textbf{عادات حماية العين}: اتبع "قاعدة 20-20-20" — كل 20 دقيقة شاشة، انظر لـ20 قدماً (نحو 6 أمتار) بعيداً 20 ثانية؛ قبل النوم بساعة لا تنظر للهاتف، الضوء الأزرق يثبط الميلاتونين ويؤثر على النوم.
  \item \textbf{إدارة الوضعية}: شاشة الهاتف بمستوى العين، لا تنحنِ؛ الاستخدام المتصل لا يتجاوز 30 دقيقة قبل النهوض للحركة.
  \item \textbf{قواعد عائلية مشتركة}: الطاولة وغرفة النوم والمكتب "مناطق بلا هاتف"؛ مشاهدة فيلم عائلي أو نشاط خارجي، تُجمع الهواتف، ليصبح "ترك الهاتف" ثقافة عائلية لا عقاباً للطفل.
\end{itemize}

\begin{tcbinfo}{الوالد أفضل "جدار حماية"}
تشير الدراسات إلى أن مدة استخدام الهاتف للوالدين مرتبطة إيجابياً بقوة مع الطفل. طلب الطفل عدم استخدام الهاتف، بينما يتصفح الوالد الفيديوهات القصيرة حتى منتصف الليل، يُفشل القواعد. لتغيير الطفل، ابدأ بترك هاتفك.

\end{tcbinfo}

\booksection{سادساً: التراث الرقمي والأطفال}

هذا القسم يُغفل بسهولة لكنه مهم. أصول الطفل الرقمية (الحسابات والصور والأرصدة والمعدات في الألعاب) تواجه قضايا الإرث والخصوصية:

\begin{itemize}
  \item \textbf{إرث الأصول الرقمية للطفل}: عند وفاة الطفل، إجراءات إرث حساباته الاجتماعية والألعاب وصور السحابة كالبالغين، يلزم شهادة وفاة وإثبات قرابة. يُنصح الوالد بضبط "جهة اتصال الإرث" أو "جهة الطوارئ" للحسابات المهمة للطفل مسبقاً (مثل جهة إرث Apple وإرث حساب WeChat).
  \item \textbf{الخصوصية الرقمية للقاصرين}: قانون حماية المعلومات الشخصية ينص على أن معالجة معلومات القاصر تحت 14 سنة تتطلب موافقة الوصي، ووضع قواعد خصوصية خاصة. راجع الوالد دورياً سياسات خصوصية تطبيقات الطفل، وأغلق "التوصية الشخصية" و"مشاركة البيانات".
  \item \textbf{نصيحة تسجيل الحساب}: الحسابات المسجلة للطفل يجب أن تستخدم بريداً/هاتف الوالد للتحكم، لتسهيل الإدارة الموحدة قبل البلوغ؛ بعد البلوغ يمكن نقل الحساب.
\end{itemize}

\begin{tcbwarn}{احذر الجانب الآخر لـ"التراث الرقمي للطفل"}
الطفل يُسجَّل منذ الولادة — صور النمو التي ينشرها الوالد وفيديوهات الروضة وبيانات ساعة الطفل الصحية. هذه البيانات يصعب حذفها تماماً بعد الرفع. قبل النشر فكّر: بعد عشر سنوات هل سيحرج الطفل من هذه الصورة؟ احترام شخصية الطفل الرقمية يبدأ بضبط “نشر صور الطفل”.

\end{tcbwarn}

\begin{tcbsummary}{خلاصة الفصل}
\begin{itemize}
  \item لا إجابة قياسية لعمر إعطاء الهاتف، لكن 3-6 تجنب، وقبل 13 سنة ساعة طفل أو هاتف وظائف انتقال
  \item "يكفي ما يكفي": الهاتف الأول ألف يوان لا رائد أعمال، الأداء الفائض يعني الإغراء الفائض
  \item iOS وقت الشاشة وAndroid Family Link/رقابة العلامة خطان دفاعيان تقنيان، استخدمهما معاً
  \item وضع الشباب ومنع إدمان الألعاب قانون، لكن يحتاج الوالد لتفعيله وحماية كلمة المرور
  \item آداب الإنترنت أهم من الحجب التقني: علّم الطفل التعرف على التنمر وحماية الخصوصية ونقد المحتوى
  \item العادات الصحية تبدأ بالوالد — تترك هاتفك، يترك الطفل هاتفه
  \item تراث الطفل الرقمي وخصوصيته قضية جديدة، اضبط "جهة إرث" مسبقاً وضبط نشر الصور
\end{itemize}

\end{tcbsummary}
